\section{Randomized --- Coins}
%%% generated by the blueprint pipeline from TCSlib.Complexity.Randomized.Coins
\subsection*{Randomized --- Coins: Introduction}

Declarations of the Lean module \texttt{TCSlib.Complexity.Randomized.Coins}.

\subsection*{Randomized --- Coins: Results}

\begin{definition}[Run of the rho coin procedure]
\lean{Randomized.rhoCoinRun}
\statementsource{arora-barak-ch07-randomized-computation}{pdf-26d95080dbde-p010-b002}
\label{Randomized.rhoCoinRun}
\leanok
Given an infinite Boolean sequence $p$ and a finite Boolean list $w$, scan $w$ against
the sequence from its beginning. If the list is empty, return no Boolean value.
Otherwise, if its first bit agrees with the current bit of $p$, discard that bit and
continue with the remaining list and the tail of $p$; if they disagree, return the
current bit of $p$.
\end{definition}

\begin{definition}[Binary value of a Boolean sequence]
\lean{Randomized.binVal}
\statementsource{arora-barak-ch07-randomized-computation}{pdf-26d95080dbde-p010-b002}
\label{Randomized.binVal}
\leanok
For an infinite Boolean sequence $p$, define its binary value by the convergent series
$\sum_{i=0}^{\infty} \mathbf{1}_{\{p_i=\mathrm{true}\}}/2^{i+1}$, where the indicator is
$1$ when the indexed bit is true and $0$ otherwise.
\end{definition}

\begin{definition}[Binary digits from a real number]
\lean{Randomized.binDigit}
\statementsource{arora-barak-ch07-randomized-computation}{pdf-26d95080dbde-p010-b002}
\label{Randomized.binDigit}
\leanok
For a real number $\rho$ and a natural number $i$, define the $i$th binary digit to be
true exactly when the remainder modulo $2$ of the greatest natural number at most
$2^{i+1}\rho$ is $1$; otherwise define it to be false.
\end{definition}

\begin{theorem}[Prefix stability of the coin run]
\lean{Randomized.rhoCoinRun_isPrefixStable}
\label{Randomized.rhoCoinRun_isPrefixStable}
\leanok
\uses{Randomized.IsPrefixStable, Randomized.rhoCoinRun}
\difficulty{4}
Let $p$ be an infinite sequence of bits. Given a finite bit string, the coin-run
procedure compares it in order with the corresponding initial bits of $p$, returning the
bit of $p$ at the first position where they differ, and returning no value if there is
no such position. This procedure is prefix-stable: any extension of a string on which it
returns a bit returns that same bit.
\end{theorem}

\begin{theorem}[Nonhalting probability for the coin run]
\lean{Randomized.notHaltedProb_rhoCoinRun}
\label{Randomized.notHaltedProb_rhoCoinRun}
\leanok
\uses{Randomized.notHaltedProb, Randomized.randProb, Randomized.randProb_congr, Randomized.randProb_false, Randomized.randProb_succ, Randomized.randProb_true, Randomized.rhoCoinRun}
\difficulty{5}
For every infinite binary sequence $p$ and every natural number $K$, the probability
that the procedure which reads a bit string and returns no output exactly when every bit
it reads agrees with the corresponding initial bit of $p$ does not halt on a uniformly
random string of $K$ bits is $(1/2)^K$.
\end{theorem}

\begin{theorem}[Expected runtime of the mismatch procedure]
\lean{Randomized.expTimeLe_rhoCoinRun}
\label{Randomized.expTimeLe_rhoCoinRun}
\leanok
\uses{Randomized.ExpTimeLe, Randomized.notHaltedProb_rhoCoinRun, Randomized.rhoCoinRun}
\difficulty{3}
Let $p:\mathbb{N}\to\{\mathrm{false},\mathrm{true}\}$ be an infinite Boolean sequence.
On a finite input string, compare its bits with the corresponding initial terms of $p$;
halt at the first mismatch and return the corresponding term of $p$, and do not halt if
there is no mismatch. For every finite cutoff $K$, the sum of the probabilities that
this procedure has not halted after $s$ bits, over $0\le s<K$, is at most $2$.
\end{theorem}

\begin{theorem}[Finite coin-run probability as a sum]
\lean{Randomized.randProb_rhoCoinRun_eq_sum}
\label{Randomized.randProb_rhoCoinRun_eq_sum}
\leanok
\uses{Randomized.randProb, Randomized.randProb_congr, Randomized.randProb_false, Randomized.randProb_succ, Randomized.randProb_true, Randomized.rhoCoinRun}
\difficulty{5}
Let $p=(p_i)_{i\ge 0}$ be a sequence of Boolean bits and let $K\ge 0$. Choose a
length-$K$ Boolean word uniformly at random, and scan it against $p$ in order, returning
the first bit $p_i$ at a position where the word differs from $p_i$. The probability
that this returned bit is true is $\displaystyle\sum_{i=0}^{K-1}\frac{\mathbf
1_{\{p_i=\mathrm{true}\}}}{2^{i+1}}$.
\end{theorem}

\begin{theorem}[Uniform bound on expected weighted nonhalting cost]
\lean{Randomized.exists_expCost_rhoCoinRun_le}
\label{Randomized.exists_expCost_rhoCoinRun_le}
\leanok
\uses{Randomized.notHaltedProb, Randomized.notHaltedProb_rhoCoinRun, Randomized.rhoCoinRun}
\difficulty{4}
For every natural number $c$, there is a rational number $C$ such that, for every
infinite bit sequence $p$ and every natural number $K$, the sum over $0\le t<K$ of
$(t+1)^c$ times the probability that the run following $p$ has not halted after $t$
random bits is at most $C$.
\end{theorem}

\begin{theorem}[Summability of binary values]
\lean{Randomized.summable_binVal}
\label{Randomized.summable_binVal}
\leanok
\difficulty{3}
For every sequence $p:\mathbb{N}\to\{\mathrm{false},\mathrm{true}\}$, the series whose
$i$th term is $2^{-(i+1)}$ if $p(i)$ is true and $0$ otherwise is summable.
\end{theorem}

\begin{theorem}[Real-valued finite coin-run probability]
\lean{Randomized.randProb_rhoCoinRun_eq_sum_real}
\label{Randomized.randProb_rhoCoinRun_eq_sum_real}
\leanok
\uses{Randomized.randProb, Randomized.randProb_rhoCoinRun_eq_sum, Randomized.rhoCoinRun}
\difficulty{3}
Let $p=(p_i)_{i\ge 0}$ be a sequence of Boolean bits and let $K\ge 0$. Choose a
length-$K$ Boolean word uniformly at random, and scan it against $p$ in order, returning
the first bit $p_i$ at a position where the word differs from $p_i$. The probability
that this returned bit is true, viewed as a real number, equals
$\displaystyle\sum_{i=0}^{K-1}\frac{\mathbf 1_{\{p_i=\mathrm{true}\}}}{2^{i+1}}$.
\end{theorem}

\begin{theorem}[Finite coin-run probability bound]
\lean{Randomized.randProb_rhoCoinRun_le_binVal}
\label{Randomized.randProb_rhoCoinRun_le_binVal}
\leanok
\uses{Randomized.binVal, Randomized.randProb, Randomized.randProb_rhoCoinRun_eq_sum_real, Randomized.rhoCoinRun, Randomized.summable_binVal}
\difficulty{2}
Let $p=(p_i)_{i\ge 0}$ be a sequence of Boolean bits, and define its binary value by
$\displaystyle\sum_{i=0}^{\infty}\mathbf 1_{\{p_i=\mathrm{true}\}}2^{-(i+1)}$. For any
$K\ge 0$, the probability that a uniformly chosen length-$K$ Boolean word, scanned
against $p$ in order, first differs at a position where the corresponding bit of $p$ is
true, is at most this binary value.
\end{theorem}

\begin{theorem}[Binary value bounded by a coin-run probability]
\lean{Randomized.binVal_le_randProb_rhoCoinRun_add}
\label{Randomized.binVal_le_randProb_rhoCoinRun_add}
\leanok
\uses{Randomized.binVal, Randomized.randProb, Randomized.randProb_rhoCoinRun_eq_sum_real, Randomized.rhoCoinRun, Randomized.summable_binVal}
\difficulty{4}
Let $p$ be an infinite Boolean sequence and let $K$ be a natural number. Define its
binary value by $\sum_{i\ge0}(1_{p(i)=\mathrm{true}})2^{-(i+1)}$. Choose a Boolean word
of length $K$ uniformly at random, and scan it against the first $K$ bits of $p$: at the
first mismatch, return the corresponding bit of $p$; if there is no mismatch, return no
value. The binary value of $p$ is at most the probability that this procedure returns
true plus $2^{-K}$.
\end{theorem}

\begin{theorem}[Limit of the matching-run probability]
\lean{Randomized.tendsto_randProb_rhoCoinRun}
\label{Randomized.tendsto_randProb_rhoCoinRun}
\leanok
\uses{Randomized.binVal, Randomized.binVal_le_randProb_rhoCoinRun_add, Randomized.randProb, Randomized.randProb_rhoCoinRun_le_binVal, Randomized.rhoCoinRun}
\difficulty{3}
Let $p:\mathbb{N}\to\{0,1\}$. For each $K$, choose uniformly a binary word of length $K$
and compare it with the first $K$ bits of $p$; the run returns true precisely when the
first differing bit of the word is $1$, and returns no value if there is no differing
bit. As $K$ tends to infinity, the probability that the run returns true converges to
the real number $\sum_{i\ge 0}p_i2^{-(i+1)}$, with bits interpreted as $0$ or $1$.
\end{theorem}

\begin{theorem}[Value of the all-ones binary sequence]
\lean{Randomized.binVal_const_true}
\label{Randomized.binVal_const_true}
\leanok
\uses{Randomized.binVal}
\difficulty{3}
The sum over all $i\ge0$ of the weights $2^{-(i+1)}$, corresponding to a sequence whose
every digit is true, equals $1$.
\end{theorem}

\begin{theorem}[Finite sum of binary digits]
\lean{Randomized.sum_binDigit}
\label{Randomized.sum_binDigit}
\leanok
\uses{Randomized.binDigit}
\difficulty{5}
Let $\rho<1$ be a real number and let $K$ be a natural number. The sum, for $0\le i<K$,
of the $i$th binary digit of $\rho$ weighted by $2^{-(i+1)}$ equals $\lfloor
2^K\rho\rfloor/2^K$, where the digit is $1$ when the parity of the integer part of
$2^{i+1}\rho$ is odd and is $0$ otherwise.
\end{theorem}

\begin{theorem}[Finite coin-run probability]
\lean{Randomized.randProb_rhoCoinRun_binDigit}
\label{Randomized.randProb_rhoCoinRun_binDigit}
\leanok
\uses{Randomized.binDigit, Randomized.randProb, Randomized.randProb_rhoCoinRun_eq_sum_real, Randomized.rhoCoinRun, Randomized.sum_binDigit}
\difficulty{3}
Let $\rho<1$ be real and let $K$ be a nonnegative integer. For each $i\ge 0$, let the
$i$th bit be true exactly when the parity of $\lfloor 2^{i+1}\rho\rfloor$ is odd. Choose
a length-$K$ Boolean word uniformly at random, and scan it against these bits in order,
returning the first scanned bit of the sequence at the first position where the word
differs from it. The probability that this returned bit is true is $\lfloor
2^K\rho\rfloor/2^K$.
\end{theorem}

\begin{theorem}[Binary expansion from floor digits]
\lean{Randomized.binVal_binDigit}
\label{Randomized.binVal_binDigit}
\leanok
\uses{Randomized.binDigit, Randomized.binVal, Randomized.binVal_le_randProb_rhoCoinRun_add, Randomized.randProb_rhoCoinRun_binDigit, Randomized.randProb_rhoCoinRun_le_binVal}
\difficulty{5}
For a real number $\rho$ with $0\le\rho<1$, define the $i$th binary digit to be true
exactly when the integer part of $2^{i+1}\rho$ is odd. The sum over all $i\ge0$ of these
digits, weighted by $2^{-(i+1)}$, equals $\rho$.
\end{theorem}

\begin{definition}[Finite biased Boolean probability]
\lean{Randomized.biasedProb}
\statementsource{arora-barak-ch07-randomized-computation}{
  pdf-26d95080dbde-p010-b002,
  pdf-26d95080dbde-p010-b003
}
\label{Randomized.biasedProb}
\leanok
For a real parameter $\rho$, define the probability of a property $A$ of Boolean lists
of length $m$ recursively. For $m=0$, it is $1$ if $A$ holds of the empty list and $0$
otherwise. For length $m+1$, it is $\rho$ times the length-$m$ probability of $A$ after
prepending true, plus $1-\rho$ times the length-$m$ probability of $A$ after prepending
false.
\end{definition}

\begin{theorem}[Recursive formula for biased probability]
\lean{Randomized.biasedProb_succ}
\label{Randomized.biasedProb_succ}
\leanok
\uses{Randomized.biasedProb}
\difficulty{1}
For any real bias parameter $\rho$, any nonnegative integer $m$, and any predicate $A$
on Boolean strings, its biased probability on strings of length $m+1$ is $\rho$ times
the biased probability of the predicate obtained by prepending true to strings of length
$m$, plus $1-\rho$ times the biased probability of the predicate obtained by prepending
false.
\end{theorem}

\begin{theorem}[Biased probability respects event equivalence]
\lean{Randomized.biasedProb_congr}
\label{Randomized.biasedProb_congr}
\leanok
\uses{Randomized.biasedProb, Randomized.biasedProb_succ}
\difficulty{4}
Let $\rho\in\bbr$, let $m$ be a natural number, and let $A,B$ be predicates on finite
Boolean words. If $A$ and $B$ agree on every word of length $m$, then their
probabilities under $m$ independent bits that are true with probability $\rho$ are
equal.
\end{theorem}

\begin{theorem}[Probability of the certain event]
\lean{Randomized.biasedProb_true}
\label{Randomized.biasedProb_true}
\leanok
\uses{Randomized.biasedProb, Randomized.biasedProb_succ}
\difficulty{4}
For any real bias parameter $\rho$ and any nonnegative integer $m$, the biased
probability of the event that is always true on Boolean strings of length $m$ is $1$.
\end{theorem}

\begin{theorem}[Probability of the empty event]
\lean{Randomized.biasedProb_false}
\label{Randomized.biasedProb_false}
\leanok
\uses{Randomized.biasedProb, Randomized.biasedProb_succ}
\difficulty{4}
For any real bias parameter $\rho$ and any nonnegative integer $m$, the biased
probability of the event that is always false on Boolean strings of length $m$ is $0$.
\end{theorem}

\begin{theorem}[Biased probability of a word]
\lean{Randomized.biasedProb_word}
\label{Randomized.biasedProb_word}
\leanok
\uses{Randomized.biasedProb, Randomized.biasedProb_congr, Randomized.biasedProb_succ}
\difficulty{5}
For a real parameter $\rho$ and any finite Boolean word $u$, its probability under the
recursively defined biased word measure is the product of its bit weights: each true bit
contributes $\rho$ and each false bit contributes $1-\rho$.
\end{theorem}

\begin{theorem}[Unbiased probability agrees with uniform probability]
\lean{Randomized.biasedProb_half}
\label{Randomized.biasedProb_half}
\leanok
\uses{Randomized.biasedProb, Randomized.biasedProb_succ, Randomized.randProb, Randomized.randProb_congr, Randomized.randProb_false, Randomized.randProb_succ, Randomized.randProb_true}
\difficulty{4}
For any nonnegative integer $m$ and any predicate $A$ on Boolean strings, the biased
probability with parameter $1/2$ of $A$ on strings of length $m$ equals the fraction of
those strings satisfying $A$.
\end{theorem}

\begin{definition}[Von Neumann extraction procedure]
\lean{Randomized.vonNeumannRun}
\statementsource{arora-barak-ch07-randomized-computation}{pdf-26d95080dbde-p010-b003}
\label{Randomized.vonNeumannRun}
\leanok
Given a finite Boolean list, read its bits in consecutive pairs. Ignore any pair whose
bits agree; at the first pair whose bits differ, return its first bit. Return no Boolean
value if there is no such pair, including when fewer than two bits remain.
\end{definition}

\begin{theorem}[Two-step biased probability expansion]
\lean{Randomized.biasedProb_two}
\label{Randomized.biasedProb_two}
\leanok
\uses{Randomized.biasedProb}
\difficulty{2}
For a real bias parameter $\rho$, let $P_m(A)$ be the recursively defined biased
probability of a predicate $A$ on finite Boolean words: the empty word has probability
$1$ or $0$ according as it satisfies $A$ or not, and for positive length the first bit
is true with weight $\rho$ and false with weight $1-\rho$. For every natural number $K$
and every such predicate $A$, the probability on words of length $2(K+1)$ is the sum
over the four possible first two bits, with their respective weights:
$P_{2(K+1)}(A)=\rho\bigl(\rho
P_{2K}(A_{11})+(1-\rho)P_{2K}(A_{10})\bigr)+(1-\rho)\bigl(\rho
P_{2K}(A_{01})+(1-\rho)P_{2K}(A_{00})\bigr)$, where $A_{bc}(w)$ means that $A$ holds of
the word obtained by prefixing $b,c$ to $w$.
\end{theorem}

\begin{theorem}[Probability of returning no bit]
\lean{Randomized.vonNeumann_none}
\label{Randomized.vonNeumann_none}
\leanok
\uses{Randomized.biasedProb, Randomized.biasedProb_false, Randomized.biasedProb_two, Randomized.vonNeumannRun}
\difficulty{4}
Let $\rho$ be any real number and let $K$ be a natural number. For $2K$ independent
bits, each equal to true with probability $\rho$, the rule that examines consecutive
pairs, discards equal pairs, and returns the first bit of the first unequal pair returns
no bit with probability $(\rho^2+(1-\rho)^2)^K$.
\end{theorem}

\begin{theorem}[Probability of returning true]
\lean{Randomized.vonNeumann_heads}
\label{Randomized.vonNeumann_heads}
\leanok
\uses{Randomized.biasedProb, Randomized.biasedProb_false, Randomized.biasedProb_true, Randomized.biasedProb_two, Randomized.vonNeumannRun}
\difficulty{4}
Let $\rho$ be any real number and let $K$ be a natural number. For $2K$ independent
bits, each equal to true with probability $\rho$, the rule that examines consecutive
pairs, discards equal pairs, and returns the first bit of the first unequal pair returns
true with probability $\frac{1-(\rho^2+(1-\rho)^2)^K}{2}$.
\end{theorem}

\begin{theorem}[Probability of returning false]
\lean{Randomized.vonNeumann_tails}
\label{Randomized.vonNeumann_tails}
\leanok
\uses{Randomized.biasedProb, Randomized.biasedProb_false, Randomized.biasedProb_true, Randomized.biasedProb_two, Randomized.vonNeumannRun}
\difficulty{4}
Let $\rho$ be any real number and let $K$ be a natural number. For $2K$ independent
bits, each equal to true with probability $\rho$, the rule that examines consecutive
pairs, discards equal pairs, and returns the first bit of the first unequal pair returns
false with probability $\frac{1-(\rho^2+(1-\rho)^2)^K}{2}$.
\end{theorem}

\begin{theorem}[Equal probabilities of either returned bit]
\lean{Randomized.vonNeumann_heads_eq_tails}
\label{Randomized.vonNeumann_heads_eq_tails}
\leanok
\uses{Randomized.biasedProb, Randomized.vonNeumannRun, Randomized.vonNeumann_heads, Randomized.vonNeumann_tails}
\difficulty{1}
Let $\rho$ be any real number and let $K$ be a natural number. For $2K$ independent
bits, each equal to true with probability $\rho$, the rule that examines consecutive
pairs, discards equal pairs, and returns the first bit of the first unequal pair returns
true and false with equal probabilities.
\end{theorem}

\begin{theorem}[Von Neumann extraction tends to fair heads]
\lean{Randomized.tendsto_vonNeumann_heads}
\label{Randomized.tendsto_vonNeumann_heads}
\leanok
\uses{Randomized.biasedProb, Randomized.vonNeumannRun, Randomized.vonNeumann_heads}
\difficulty{4}
Let $0<\rho<1$. For each $K$, draw $2K$ independent bits, each equal to $1$ with
probability $\rho$, and inspect successive pairs, discarding equal pairs and returning
the first bit of the first unequal pair; if there is no unequal pair, return no value.
As $K$ tends to infinity, the probability that this procedure returns $1$ converges to
$1/2$.
\end{theorem}

\begin{theorem}[Expected pairs before rejection]
\lean{Randomized.vonNeumann_expected_pairs}
\label{Randomized.vonNeumann_expected_pairs}
\leanok
\uses{Randomized.biasedProb, Randomized.vonNeumannRun, Randomized.vonNeumann_none}
\difficulty{4}
Let $0<\rho<1$, and let $M$ be a natural number. For each $K<M$, consider a sequence of
$2K$ independent bits, each equal to true with probability $\rho$, and apply the rule
that examines consecutive pairs, discards equal pairs, and stops at the first unequal
pair, returning its first bit; the result is none if no unequal pair occurs. The sum,
over $K=0,\ldots,M-1$, of the probabilities that this rule returns none is at most
$\frac{1}{2\rho(1-\rho)}$.
\end{theorem}

\begin{theorem}[Uniform Bound for an Exponentially Weighted Power Sum]
\lean{Randomized.exists_sum_rpow_div_two_pow_le}
\label{Randomized.exists_sum_rpow_div_two_pow_le}
\leanok
\difficulty{4}
For every real number $c$, there is a real constant $D$ such that for every natural
number $n \ge 1$, $\sum_{i=1}^{n} i^c / 2^i \le D$.
\end{theorem}
